\chapter{Basit Doğrusal Regresyon}
\index{doğrusal regresyon}
\index{basit doğrusal regresyon}

\begin{hedefler}
Bu bölümün sonunda okuyucu:
\begin{itemize}
  \item basit doğrusal regresyonda yanıt, açıklayıcı değişken, katsayı ve hata terimini ayırt edebilecek,
  \item kesme ve eğim katsayılarını değişkenlerin özgün birimleriyle yorumlayabilecek,
  \item en küçük kareler doğrusunun nasıl seçildiğini ve artıkların rolünü açıklayabilecek,
  \item eğim için standart hata, hipotez testi ve güven aralığını birlikte değerlendirebilecek,
  \item belirlilik katsayısı ile artık standart hatasını model uyumunun farklı yönleri olarak yorumlayabilecek,
  \item ortalama yanıt güven aralığı ile bireysel tahmin aralığını ayırabilecek,
  \item doğrusallık, sabit varyans, bağımsızlık, normallik ve etkili gözlem koşullarını tanılayabilecek,
  \item öngörü, ilişki ve nedensel etki kavramlarını birbirine karıştırmadan Python ve SPSS sonuçlarını raporlayabilecektir.
\end{itemize}
\end{hedefler}

\begin{hazirlik}
Bir regresyon doğrusunun bütün gözlem noktalarından geçmesi gerekir mi?
``En iyi uyum'' ifadesinin hangi ölçüte göre tanımlanabileceğini ve bazı
noktaların neden doğrunun üstünde ya da altında kalacağını düşünün.
\end{hazirlik}

\section{PDR vakası: Uyku süresinden kaygı puanını yordamlamak}

Önceki bölümde 20 öğrencinin gecelik uyku süresi ile akademik kaygı puanı
arasında $r=-0{,}585$ korelasyon bulunmuştu. Korelasyon ilişkinin yönünü ve
standartlaştırılmış derecesini özetledi. Psikolojik danışma birimi şimdi daha
somut sorular sormaktadır:

\begin{enumerate}
  \item Uyku süresi bir saat arttığında beklenen kaygı puanı ortalama kaç puan değişir?
  \item Yedi saat uyuyan öğrencilerin ortalama kaygı puanı nedir?
  \item Yedi saat uyuyan tek bir öğrencinin puanı hangi geniş aralıkta olabilir?
  \item Doğrusal model veriyi yeterince iyi temsil ediyor mu?
\end{enumerate}


Basit doğrusal regresyon bu soruları yanıt ve açıklayıcı değişken rollerini
belirleyerek ele alır. Bu örnekte akademik kaygı yanıt değişkeni $Y$, uyku
süresi açıklayıcı değişken $X$ olacaktır.

\begin{researchbox}
``Uyku süresinden kaygıyı yordamlamak'' ifadesi nedensel etki iddiası değildir.
Kesitsel gözlemsel veride regresyon eğimi, uyku süresi farklı öğrencilerin
ortalama kaygı farkını özetler. Aynı öğrencinin uykusu bir saat artırıldığında
kaygısının bu miktarda değişeceği sonucuna doğrudan varılamaz.
\end{researchbox}

\section{İlişkiden doğrusal modele}
\index{yanıt değişkeni}
\index{açıklayıcı değişken}
\index{hata terimi}

Basit doğrusal regresyon, $X=x$ olan birimlerin ortalama yanıtını
\[
E(Y\mid X=x)=\beta_0+\beta_1x
\]
biçiminde modeller. Tek tek gözlemler için
\[
Y_i=\beta_0+\beta_1X_i+\varepsilon_i
\]
yazılır. Bu modelin tahmin ve çıkarım çerçevesi doğrusal model kaynaklarında
ayrıntılı olarak ele alınır \cite{kutner2005}. Modelin bileşenleri şunlardır:

\begin{description}
  \item[$Y_i$] $i$'inci öğrencinin gözlenen yanıtı, burada kaygı puanı,
  \item[$X_i$] açıklayıcı değişken değeri, burada uyku süresi,
  \item[$\beta_0$] evrendeki regresyon doğrusunun kesme katsayısı,
  \item[$\beta_1$] evrendeki regresyon eğimi,
  \item[$\varepsilon_i$] öğrencinin koşullu ortalamadan ayrılan kısmı.
\end{description}

Hata terimi yalnız ölçüm hatası değildir. Modelde bulunmayan bireysel
özellikler, gün içi değişkenlik ve doğrusal ortalama çevresindeki doğal
farklılıklar da hata teriminde yer alır.

Örneklemden tahmin edilen doğru
\[
\widehat{Y}=b_0+b_1X
\]
biçimindedir. Yirmi öğrencilik uyku--kaygı verisi için
\[
\widehat{\text{Kaygı}}=89{,}54-4{,}98\times\text{Uyku}
\]
elde edilir.

\section{Kesme ve eğim katsayısını yorumlamak}
\index{regresyon!kesme katsayısı}
\index{regresyon!eğim katsayısı}
\index{merkezleme}

\subsection{Eğim}

$b_1=-4{,}98$, uyku süresi bir saat daha yüksek olan öğrencilerin beklenen
kaygı puanının örneklemde ortalama 4,98 puan daha düşük olduğunu gösterir.
Yorum hem X'in hem Y'nin birimini içermelidir.

Eğim ``her öğrenci bir saat fazla uyursa kaygısı kesin 4,98 puan azalır''
anlamına gelmez. Regresyon doğrusu koşullu ortalamayı tanımlar; bireysel
puanlar doğrunun çevresinde dağılır. Gözlemsel tasarımda eğim aynı zamanda
nedensel değişim olarak yorumlanamaz.

\subsection{Kesme}

$b_0=89{,}54$, modelde uyku süresi 0 saat olduğunda beklenen kaygı puanıdır.
Ancak veride uyku süresi 4,5 ile 8,5 saat arasındadır. Sıfır saat veri
aralığının çok dışında ve uygulamada olağandışı olduğundan kesmenin doğrudan
PDR yorumu anlamlı değildir. Kesme yine de doğrunun matematiksel konumunu
belirler.

X değişkenini ortalaması çevresinde merkezlemek yorumu kolaylaştırabilir.
$X_c=X-6{,}67$ tanımlanırsa yeni kesme, ortalama uyku süresindeki beklenen
kaygı puanı olan 56,3'e eşit olur; eğim değişmez.

\begin{mistakebox}
Kesme katsayısını yalnız yazılım verdiği için yorumlamayın. $X=0$ değeri
ölçekte olanaksız, veri aralığının dışında veya araştırma açısından anlamsızsa
kesme yalnız modelin cebirsel bileşeni olarak belirtilmelidir.
\end{mistakebox}

\section{En küçük kareler doğrusu}
\index{en küçük kareler yöntemi}
\index{artık}
\index{uydurulan değer}

Her öğrenci için uydurulan değer
\[
\hat Y_i=b_0+b_1X_i
\]
ve artık
\[
e_i=Y_i-\hat Y_i
\]
olarak tanımlanır. Pozitif artık gözlenen puanın model tahmininden yüksek,
negatif artık düşük olduğunu gösterir.

\begin{figure}[htbp]
\centering
\begin{tikzpicture}[alt={Bir gözlemin regresyon doğrusu üzerindeki tahmini ile gözlenen değeri arasındaki artığı gösteren çizim},x=1.0cm,y=0.12cm]
  \draw[->,PrismGray] (4,35)--(9,35) node[right,font=\scriptsize]{uyku};
  \draw[->,PrismGray] (4,35)--(4,78) node[above,font=\scriptsize]{kaygı};
  \draw[PrismBlue,very thick] (4.5,67.1)--(8.5,47.2);
  \fill[PrismWarnOrange!85!black] (7,62) circle (2.4pt);
  \fill[PrismBlue] (7,54.7) circle (2.2pt);
  \draw[<->,PrismWarnOrange!90!black,thick] (7,54.7)--(7,62);
  \node[right,font=\scriptsize,text=PrismWarnOrange!80!black] at (7.08,58.4) {artık $e_i$};
  \node[below,font=\scriptsize,text=PrismBlue] at (7,54.2) {$\hat Y_i$};
\end{tikzpicture}
\caption{Bir gözlem için gözlenen değer, uydurulan değer ve artık.}
\end{figure}

En küçük kareler yöntemi
\[
SSE=\sum_{i=1}^ne_i^2
=\sum_{i=1}^n(Y_i-b_0-b_1X_i)^2
\]
toplamını en aza indiren katsayıları seçer. Artıkların karelenmesi pozitif ve
negatif sapmaların birbirini götürmesini engeller ve büyük sapmalara daha çok
ağırlık verir.

Çözüm
\[
b_1=\frac{\sum(X_i-\bar X)(Y_i-\bar Y)}
{\sum(X_i-\bar X)^2}
=r\frac{s_Y}{s_X},
\qquad
b_0=\bar Y-b_1\bar X
\]
biçimindedir. Uydurulan doğru $(\bar X,\bar Y)$ noktasından geçer.

Kesme terimli en küçük kareler modelinde artıkların toplamı sıfırdır ve
uydurulan değerlerin ortalaması gözlenen Y ortalamasına eşittir. Bu cebirsel
özellikler modelin doğru veya nedensel olduğunu garanti etmez.

\begin{sezgibox}
En küçük kareler doğrusu bütün noktalardan geçen bir çizgi aramaz. Noktaların
doğruya dikey uzaklıklarının kareleri toplamını en küçük yapan denge çizgisini
seçer. Hangi değişkenin Y olduğu bu nedenle önemlidir.
\end{sezgibox}

\section{Korelasyon ile regresyon arasındaki bağlantı}
\index{korelasyon!regresyon ilişkisi}

Tek açıklayıcılı, kesme terimli regresyonda:

\begin{itemize}
  \item eğim ile Pearson korelasyonu aynı işarete sahiptir,
  \item $b_1=r(s_Y/s_X)$ bağıntısı geçerlidir,
  \item $R^2=r^2$ olur,
  \item eğim için $H_0:\beta_1=0$ testi ile $H_0:\rho=0$ testi aynı $t$ ve $p$ sonucunu verir.
\end{itemize}

Yine de iki yöntem aynı değildir. Korelasyon X ve Y'ye simetrik davranır;
değişkenler yer değiştirince $r$ değişmez. Regresyon Y'yi X'ten yordamayı
amaçlar; değişkenler yer değiştirince eğim, artıklar ve tahminler değişir.

\section{Eğim için çıkarım}
\index{regresyon!eğim testi}
\index{regresyon!eğim güven aralığı}

Örneklem eğimi $b_1$, evren eğimi $\beta_1$ için nokta tahminidir. Sıfır
doğrusal eğim hipotezi
\[
H_0:\beta_1=0,
\qquad
H_1:\beta_1\neq0
\]
ve test istatistiği
\[
t=\frac{b_1}{SE(b_1)},
\qquad sd=n-2
\]
biçimindedir.

Uyku--kaygı modelinde
\[
b_1=-4{,}984,
\qquad SE(b_1)=1{,}629,
\]
dolayısıyla
\[
t=\frac{-4{,}984}{1{,}629}=-3{,}059,
\qquad p=0{,}007.
\]

Eğim için yüzde 95 güven aralığı
\[
b_1\pm t^*SE(b_1)
=[-8{,}406;\ -1{,}561]
\]
olarak bulunur. Veri, uyku süresindeki bir saatlik artış başına beklenen kaygı
puanında yaklaşık 1,6 ile 8,4 puanlık azalmalarla uyumludur. Aralığın genişliği
etkinin büyüklüğü hakkındaki belirsizliği gösterir.

\begin{mistakebox}
$p=0{,}007$ eğimin doğru olma olasılığının yüzde 99,3 olduğu anlamına gelmez.
Sıfır eğim ve model varsayımları altında gözlenen kadar uç bir test
istatistiğinin olasılığını özetler.
\end{mistakebox}

\section{Değişkenliğin ayrıştırılması ve model uyumu}
\index{R kare@$R^2$}
\index{model uyumu}
\index{SSE}
\index{SSR}
\index{SST}

Yanıt değişkenindeki toplam değişkenlik
\[
SST=\sum(Y_i-\bar Y)^2
\]
olup model tarafından temsil edilen ve artıkta kalan parçalara ayrılır:
\[
SST=SSR+SSE.
\]

Uyku--kaygı verisinde
\[
SST=1722{,}20,
\qquad SSR=589{,}16,
\qquad SSE=1133{,}04.
\]

Belirlilik katsayısı
\[
R^2=\frac{SSR}{SST}
=1-\frac{SSE}{SST}
\]
olduğundan
\[
R^2=0{,}342
\]
bulunur. Örneklemde kaygı puanı değişkenliğinin yüzde 34,2'si uyku süresiyle
kurulan doğrusal model tarafından temsil edilmektedir.

$R^2$ modelin nedensel doğruluğunu, bireysel tahminlerin kesinliğini veya yeni
örneklem performansını garanti etmez. Gereksiz değişken eklemek çoklu
regresyonda $R^2$ değerini düşürmez; bu nedenle düzeltilmiş $R^2$ ve dış
örneklem doğrulaması da önemlidir.

\subsection{Artık standart hatası}
\index{artık standart hatası}

Hata varyansının tahmini
\[
MSE=\frac{SSE}{n-2}
\]
ve artık standart hatası
\[
s_e=\sqrt{MSE}
\]
ile verilir. Örnekte
\[
s_e=\sqrt{1133{,}04/18}\approx7{,}93
\]
puandır. Bu değer gözlenen kaygı puanlarının regresyon doğrusu çevresindeki
tipik dikey saçılmasını Y'nin özgün biriminde özetler.

\begin{researchbox}
Yüzde 34 açıklanan değişkenlik önemli görünebilir; fakat yaklaşık 7,9 puanlık
artık standart hatası bireysel öğrenciler arasında geniş farklılık kaldığını
gösterir. Grup düzeyindeki ilişki tek bir öğrenci için kesin puan üretmez.
\end{researchbox}

\section{Genel F testi}
\index{F testi@$F$ testi!regresyon}

Regresyon ANOVA tablosunda model ortalama karesi hata ortalama karesine
bölünür:
\[
F=\frac{MSR}{MSE}.
\]
Tek açıklayıcılı modelde eğim testinin karesi genel $F$ testine eşittir:
\[
F=t^2.
\]
Örnekte $(-3{,}059)^2=9{,}36$ ve $F(1,18)=9{,}36$, $p=0{,}007$ elde edilir.
Birden fazla açıklayıcılı modelde genel $F$ testi bütün eğimlerin birlikte
sıfır olup olmadığını değerlendirir.

\section{Uydurulan değer, güven aralığı ve tahmin aralığı}
\index{güven aralığı!ortalama yanıt}
\index{tahmin aralığı!bireysel}

Yedi saat uyku için uydurulan ortalama kaygı puanı
\[
\hat Y=89{,}54-4{,}98(7)=54{,}66
\]
olur. Bu nokta tahmini iki farklı belirsizlik sorusuna yol açar.

\subsection{Ortalama yanıt için güven aralığı}

Yedi saat uyuyan bütün benzer öğrencilerin ortalama kaygı puanı için yüzde 95
güven aralığı yaklaşık
\[
[50{,}76;\ 58{,}55]
\]
olarak bulunur. Bu aralık regresyon doğrusunun $X=7$ noktasındaki ortalama
konum belirsizliğini gösterir.

\subsection{Yeni birey için tahmin aralığı}

Yedi saat uyuyan yeni bir öğrencinin puanı için yüzde 95 tahmin aralığı
yaklaşık
\[
[37{,}54;\ 71{,}77]
\]
olur. Bireysel öğrencilerin ortalama çizgi çevresindeki doğal değişkenliğini de
içerdiği için tahmin aralığı her zaman ortalama güven aralığından geniştir.

\begin{table}[htbp]
\centering
\caption{Yedi saat uyku için iki farklı aralığın yorumu.}
\begin{tabular}{@{}p{0.27\textwidth}p{0.24\textwidth}p{0.30\textwidth}@{}}
\toprule
Aralık & Yaklaşık sınırlar & Hedef \\
\midrule
Ortalama güven aralığı & 50,76--58,55 & Yedi saat uyuyanların evren ortalaması \\
Bireysel tahmin aralığı & 37,54--71,77 & Yedi saat uyuyan yeni bir öğrenci \\
\bottomrule
\end{tabular}
\end{table}

Her iki aralık da $X$ değeri örneklem ortalamasına yakınken daha dardır; veri
aralığının uçlarına gidildikçe genişler. Tahmin aralığı model hatasını,
ortalama aralığı ise yalnız koşullu ortalamanın tahmin belirsizliğini içerir.

\begin{mistakebox}
Regresyon doğrusu üzerindeki 54,66 değeri yedi saat uyuyan her öğrencinin
puanı değildir. Aynı uyku süresine sahip öğrenciler çok farklı kaygı puanları
alabilir; bireysel tahmin aralığı bu değişkenliği görünür kılar.
\end{mistakebox}

\section{Ekstrapolasyon}
\index{ekstrapolasyon}

Model 4,5 ile 8,5 saat arasındaki veriden kurulmuştur. Dokuz veya on saat için
doğruyu uzatmak ekstrapolasyondur. İlişkinin veri aralığı dışında aynı eğimle
süreceği bilinmez. Daha uç değerlerde biyolojik sınırlar, doygunluk veya
farklı öğrenci grupları devreye girebilir.

Ekstrapolasyon yalnız X'in gözlenen minimum ve maksimumuna değil, veri
yoğunluğuna göre de değerlendirilmelidir. Sekiz saat çevresinde çok az gözlem
varsa bu bölgedeki tahmin, aralık içinde olsa bile daha belirsiz olabilir.

\begin{mistakebox}
Yüksek $R^2$ ekstrapolasyonu güvenli hâle getirmez. Model gözlenen aralıkta iyi
uyabilir, ancak aralığın dışında ilişki biçimi tamamen değişebilir.
\end{mistakebox}

\section{Artıklarla model kontrolü}
\index{artık analizi}
\index{artık grafiği}

Regresyon sonuçları yorumlanmadan önce artıklar uydurulan değerler ve X ile
birlikte incelenir. İyi davranan bir artık grafiğinde noktalar sıfır çizgisi
çevresinde örüntüsüz ve yaklaşık sabit genişlikte dağılır.

\begin{table}[htbp]
\centering
\caption{Artık grafiklerinde görülebilecek örüntüler.}
\begin{tabular}{@{}p{0.24\textwidth}p{0.28\textwidth}p{0.29\textwidth}@{}}
\toprule
Örüntü & Olası sorun & İncelenecek yaklaşım \\
\midrule
Eğri desen & Doğrusal olmayan ortalama ilişkisi & Eğrisel terim veya farklı model \\
Huni biçimi & Değişen hata varyansı & Dönüşüm veya sağlam standart hata \\
Tek büyük artık & Y yönünde olağandışı gözlem & Kayıt ve bağlam kontrolü \\
Uç X değeri & Yüksek kaldıraç & Etki ve duyarlılık incelemesi \\
Sıralı dalga & Bağımlı veya zamansal hata & Bağımlılığı modelleme \\
\bottomrule
\end{tabular}
\end{table}

Artık grafiği yalnız ``varsayımlar sağlandı/sağlanmadı'' kutusu değildir.
Modelin veriyi hangi yönden eksik temsil ettiğini gösteren tanı aracıdır.

\section{Regresyon varsayımları}
\index{regresyon!varsayımlar}
\index{doğrusallık}
\index{sabit varyans}
\index{hata normalliği}
\index{bağımsızlık}

Basit doğrusal regresyonda farklı amaçlar için farklı koşullar önemlidir:

\begin{description}
  \item[Doğrusallık] $E(Y\mid X)$ veri aralığında yaklaşık doğrusaldır.
  \item[Bağımsızlık] Birimlerin hata terimleri birbirinden bağımsızdır veya bağımlılık modellenmiştir.
  \item[Sabit varyans] Hata saçılması X boyunca yaklaşık sabittir.
  \item[Yaklaşık normallik] Küçük örneklemde klasik test ve aralıklar için hatalar yaklaşık normaldir.
  \item[Uygun ölçüm] X ve Y ölçümleri anlamlı, karşılaştırılabilir ve yeterince güvenilirdir.
\end{description}

Normallik X değişkeninin veya Y'nin tek başına normal olması anlamına gelmez;
koşullu hata dağılımına ilişkindir. Büyük örneklemde eğim çıkarımı orta düzey
normallik sapmalarına daha dayanıklı olabilir, ancak doğrusallık ve bağımsızlık
ihlalleri hâlâ önemlidir.

\begin{assumptionbox}
Artık grafiğini, Q--Q grafiğini ve etkili gözlem ölçülerini birlikte inceleyin.
Tek bir normallik veya değişen varyans testinin $p$-değeri modelin bütün
geçerliğini belirlemez. Grafik, tasarım ve araştırma hedefi birlikte
değerlendirilmelidir.
\end{assumptionbox}

\section{Aykırı, yüksek kaldıraçlı ve etkili gözlemler}
\index{aykırı değer}
\index{kaldıraç}
\index{etkili gözlem}
\index{Cook uzaklığı}

Üç kavram ayrılmalıdır:

\begin{description}
  \item[Aykırı gözlem] Y değeri modelin tahmininden uzaktır; büyük artığa sahiptir.
  \item[Yüksek kaldıraç] X değeri örneklemin X merkezinden uzaktır.
  \item[Etkili gözlem] Çıkarıldığında katsayıları veya tahminleri belirgin değiştirir.
\end{description}

Yüksek kaldıraçlı bir nokta doğru üzerinde bulunursa büyük artığa sahip
olmayabilir; yine de eğimi güçlü biçimde sabitleyebilir. Cook uzaklığı, kaldıraç
ve öğrencilendirilmiş artıklar birlikte incelenebilir. Bunlar otomatik silme
kuralları değil, araştırma sinyalleridir.

\begin{enumerate}
  \item Veri girişi, birim ve kodlama doğrulanır.
  \item Gözlemin hedef evrene ait olup olmadığı incelenir.
  \item Gözlem dahil ve hariç katsayılar ile aralıklar karşılaştırılır.
  \item Karar sonuçtan bağımsız bilimsel gerekçeyle verilir.
  \item Duyarlılık analizi şeffaf biçimde raporlanır.
\end{enumerate}

\begin{researchbox}
Olağandışı bir öğrenci, istatistiksel modeli etkilediği kadar hizmet açısından
da önemli olabilir. Gerçek bir gözlemi modeli daha ``temiz'' göstermek için
silmek, hem bilimsel bilgiyi hem danışmanlık gereksinimini görünmez kılabilir.
\end{researchbox}

\section{Değişen varyans ve sağlam standart hatalar}
\index{değişen varyans}
\index{sağlam standart hata}

Kaygı puanlarının saçılması düşük uyku düzeylerinde daha büyük, yüksek uyku
düzeylerinde daha küçük olabilir. Böyle bir huni biçiminde en küçük kareler
eğimi belirli koşullarda doğrusal ortalama için kullanılabilir; fakat klasik
standart hata ve güven aralığı yanlış olabilir.

Olası yaklaşımlar:

\begin{itemize}
  \item veri ve değişken tanımlarını yeniden kontrol etmek,
  \item kuramsal olarak uygun dönüşüm kullanmak,
  \item değişen varyansa dayanıklı standart hatalar hesaplamak,
  \item varyans yapısını doğrudan modellemek,
  \item sonucu duyarlılık analizleriyle karşılaştırmak.
\end{itemize}

Sağlam standart hata doğrusal olmayan ortalama ilişkisini düzeltmez. Önce
model biçimi, sonra çıkarım yöntemi değerlendirilmelidir.

\section{Doğrusal olmayan ilişki}
\index{doğrusal olmayan ilişki}
\index{polinom terimi}

Artık grafiğinde eğri desen varsa düz doğru ortalama ilişkiyi eksik temsil
ediyor olabilir. Örneğin uyku ile kaygı ilişkisi bir noktadan sonra
doygunlaşabilir. Kuramsal gerekçe varsa
\[
E(Y\mid X)=\beta_0+\beta_1X+\beta_2X^2
\]
gibi eğrisel model düşünülebilir.

Karesel terim eklendiğinde $\beta_1$ artık tek başına sabit eğim değildir;
X'in etkisi düzeye göre değişir. Model sonuç görüldükten sonra sınırsız terim
denenerek seçilmemeli, yeni biçim mümkünse bağımsız veride doğrulanmalıdır.

\section{Eksik veri ve ölçüm hatası}
\index{eksik veri}
\index{ölçüm hatası}

Basit regresyon çoğunlukla X veya Y'si eksik satırları analiz dışı bırakır.
Eksiklik kaygı düzeyi, uyku sorunu veya yardım arama davranışıyla ilişkiliyse
tam durum analizi yanlı olabilir. Analize giren ve girmeyen öğrencilerin
sayısı ile bilinen özellikleri karşılaştırılmalıdır.

X değişkenindeki rastgele ölçüm hatası eğimi çoğunlukla sıfıra doğru
zayıflatır. Öz-bildirim uyku süresinin hatalı olması gerçek ilişkinin daha
küçük görünmesine yol açabilir. Y'deki ölçüm hatası ise artık varyansını ve
tahmin aralıklarını genişletebilir.

\begin{mistakebox}
Daha fazla gözlem sistematik ölçüm hatasını düzeltmez. Büyük örneklem hatalı
ölçüme dayalı eğimi çok hassas biçimde tahmin edebilir.
\end{mistakebox}

\section{Öngörü ile açıklama ve nedenselliği ayırmak}
\index{öngörü}
\index{açıklama}
\index{nedensel çıkarım}

Regresyon en az üç farklı amaçla kullanılabilir:

\begin{description}
  \item[Betimleme] Örneklemdeki ortalama doğrusal ilişkiyi özetlemek,
  \item[Öngörü] Yeni birimlerin Y değerini olabildiğince doğru tahmin etmek,
  \item[Nedensel çıkarım] X'e müdahalenin Y üzerindeki etkisini tahmin etmek.
\end{description}

İyi öngören bir değişken nedensel olmayabilir; nedensel bir etki de bireysel
öngörüye az katkı sağlayabilir. Nedensel yorum için zaman sırası, karıştırıcı
değişkenler, seçim süreci ve araştırma tasarımı hakkında ek varsayımlar gerekir.

Uyku ile kaygı ilişkisinde sınav yükü her ikisini etkileyebilir. Sınav yükünü
modele eklemek yararlı olabilir; ancak ölçülmemiş karıştırıcılar, ters
nedensellik ve hatalı kovaryat ölçümü devam eder. ``Diğer değişkenler kontrol
edildi'' ifadesi nedenselliği otomatik olarak kurmaz.

\section{Modelin yeni veride değerlendirilmesi}
\index{çapraz doğrulama}
\index{eğitim verisi}
\index{test verisi}

Modelin kurulduğu örneklemdeki $R^2$ ve hata, aynı veriye uyumu gösterir. Yeni
öğrencilerde öngörü başarısı daha düşük olabilir. Öngörü amacı varsa veri
eğitim ve test bölümlerine ayrılabilir veya çapraz doğrulama kullanılabilir.

Yeni veride ortalama mutlak hata ve hata kareleri ortalamasının karekökü gibi
ölçüler incelenebilir. Çok küçük örneklemde tek bir eğitim--test bölmesi
kararsızdır; çapraz doğrulama daha verimli olabilir. Model seçimi ve başarı
ölçütü veri analizinden önce belirlenmelidir \cite{hastie2009,james2021}.

\begin{sezgibox}
Bir doğrunun geçmiş veriye çok iyi uyması, yeni öğrencileri iyi tahmin edeceği
anlamına gelmez. Modelin gerçek sınavı, kurulurken görmediği verideki
performansıdır.
\end{sezgibox}

\section{Python ile basit doğrusal regresyon}

\begin{lstlisting}
import numpy as np
import matplotlib.pyplot as plt
import statsmodels.api as sm

uyku = np.array([4.5,5.0,5.2,5.5,5.8,6.0,6.1,6.3,6.5,6.7,
                 6.8,7.0,7.1,7.3,7.5,7.6,7.8,8.0,8.2,8.5])
kaygi = np.array([68,52,75,60,66,48,70,55,64,58,
                  45,62,53,46,60,43,55,49,57,40])

X = sm.add_constant(uyku)
model = sm.OLS(kaygi, X).fit()
print(model.summary())
print(model.conf_int())

# Yedi saat uyku icin ortalama ve bireysel araliklar
yeni = np.array([[1, 7.0]])
print(model.get_prediction(yeni).summary_frame(alpha=0.05))

# Temel tani grafikleri
fig, ax = plt.subplots(1, 2, figsize=(8, 3))
ax[0].scatter(uyku, kaygi)
ax[0].plot(uyku, model.fittedvalues, color="navy")
ax[0].set_xlabel("Uyku suresi")
ax[0].set_ylabel("Kaygi puani")
ax[1].scatter(model.fittedvalues, model.resid)
ax[1].axhline(0, color="black", linestyle="--")
ax[1].set_xlabel("Uydurulan deger")
ax[1].set_ylabel("Artik")
plt.tight_layout()
plt.show()

# Etkili gozlem olculeri
etki = model.get_influence().summary_frame()
print(etki[["standard_resid", "hat_diag", "cooks_d"]])
\end{lstlisting}

\texttt{summary\_frame} çıktısında \texttt{mean\_ci\_lower} ve
\texttt{mean\_ci\_upper} ortalama yanıt aralığını;
\texttt{obs\_ci\_lower} ve \texttt{obs\_ci\_upper} yeni birey tahmin
aralığını verir. İki aralık birbirinin yerine raporlanmamalıdır.

\section{SPSS ile doğrusal regresyon}

\begin{enumerate}
  \item Önce \textbf{Graphs $>$ Chart Builder} ile saçılım grafiği oluşturun.
  \item \textbf{Analyze $>$ Regression $>$ Linear} yolunu izleyin.
  \item Kaygı puanını \textbf{Dependent}, uyku süresini \textbf{Independent(s)} alanına taşıyın.
  \item \textbf{Statistics} bölümünde tahminler, model uyumu ve güven aralıklarını seçin.
  \item \textbf{Plots} bölümünde standartlaştırılmış artıklarla uydurulan değer grafiği ve normal P--P grafiği oluşturun.
  \item \textbf{Save} bölümünde uydurulan değerler, artıklar, Cook uzaklığı ve kaldıraç ölçülerini gerektiğinde kaydedin.
  \item Model Summary, ANOVA ve Coefficients tablolarını tanı grafikleriyle birlikte okuyun.
\end{enumerate}

\begin{outputguidebox}
Çıktıyı şu sırayla okuyun: (1) örneklem büyüklüğü ve değişken kodları, (2)
saçılım grafiği, (3) Model Summary tablosundaki $R^2$ ve tahmin hatası, (4)
ANOVA tablosundaki genel $F$, (5) Coefficients tablosundaki $B$, standart hata,
$t$, $p$ ve güven aralığı, (6) artık ve etki tanıları. Standartlaştırılmamış
$B$ eğimini, standartlaştırılmış Beta katsayısıyla karıştırmayın.
\end{outputguidebox}

\section{Bütünleştirici sonuç yorumu}

Uyku--kaygı modelinin temel sonuçları şöyledir:

\begin{itemize}
  \item Denklem: $\widehat{\text{Kaygı}}=89{,}54-4{,}98\text{Uyku}$.
  \item Eğim: saat başına yaklaşık 4,98 puan daha düşük beklenen kaygı.
  \item Eğim güven aralığı: $[-8{,}41;-1{,}56]$.
  \item Eğim testi: $t(18)=-3{,}06$, $p=0{,}007$.
  \item Model uyumu: $R^2=0{,}342$, artık standart hatası 7,93 puan.
  \item Yedi saat için ortalama tahmin: 54,66, yüzde 95 GA $[50{,}76;58{,}55]$.
  \item Yedi saat için bireysel tahmin aralığı: $[37{,}54;71{,}77]$.
\end{itemize}

Bu sonuçlar örneklemde negatif doğrusal ilişki bulunduğunu ve ortalama
doğrunun sıfır eğimden farklı olduğuna ilişkin kanıtı gösterir. Geniş bireysel
tahmin aralığı yalnız uyku süresinin tek öğrencinin kaygısını kesin biçimde
yordamadığını ortaya koyar.

\section{Bilimsel raporlama örneği}

``Gecelik uyku süresi akademik kaygı puanını negatif yönde yordamıştır,
$b=-4{,}98$, $SH=1{,}63$, $t(18)=-3{,}06$, $p=0{,}007$, yüzde 95 GA
$[-8{,}41;-1{,}56]$. Model kaygı puanındaki örneklem değişkenliğinin yüzde
34,2'sini temsil etmiştir, $R^2=0{,}342$, $F(1,18)=9{,}36$.
Artık standart hatası 7,93 puandır. Artık grafikleri belirgin eğrilik
göstermemiştir. Kesitsel öz-bildirim tasarımı nedeniyle eğim nedensel uyku
etkisi olarak yorumlanmamıştır.''

Bir öngörü raporunda bunlara ek olarak yeni veride doğrulama yöntemi ve hata
ölçüsü; nedensel araştırmada ise atama, karıştırıcı seçimi ve duyarlılık
varsayımları belirtilmelidir.

\section{Bir regresyon araştırmasını değerlendirirken}

\begin{enumerate}
  \item Y ve X değişkenlerinin rolleri araştırma sorusuyla uyumlu mu?
  \item Saçılım grafiği doğrusal ilişkiyi ve veri aralığını destekliyor mu?
  \item Eğim özgün birimleri ve güven aralığıyla yorumlanmış mı?
  \item Kesme yalnız $X=0$ anlamlıysa mı yorumlanmış?
  \item $R^2$ nedensellik veya bireysel tahmin kesinliği gibi sunulmuş mu?
  \item Artık, değişen varyans ve etkili gözlem tanıları incelenmiş mi?
  \item Ortalama güven aralığı ile bireysel tahmin aralığı ayrılmış mı?
  \item Ekstrapolasyon veya yeni veriye doğrulanmamış öngörü yapılmış mı?
  \item İlişki, öngörü ve nedensellik amaçları birbirinden ayrılmış mı?
\end{enumerate}

\section{R ile basit doğrusal regresyon}
\label{sec:r-regression}

\textbf{Soru.} Uyku süresinden kaygı puanını yordayan doğrusal modeli kurun;
katsayı aralıklarını, yedi saat için ortalama güven aralığını, bireysel tahmin
aralığını ve temel tanı grafiklerini üretin.

\begin{lstlisting}[style=prismR]
veri <- data.frame(
  uyku = c(4.5,5.0,5.2,5.5,5.8,6.0,6.1,6.3,6.5,6.7,
           6.8,7.0,7.1,7.3,7.5,7.6,7.8,8.0,8.2,8.5),
  kaygi = c(68,52,75,60,66,48,70,55,64,58,
            45,62,53,46,60,43,55,49,57,40)
)
model <- lm(kaygi ~ uyku, data = veri)
summary(model)
confint(model)

yeni <- data.frame(uyku = 7)
predict(model, newdata = yeni, interval = "confidence")
predict(model, newdata = yeni, interval = "prediction")

par(mfrow = c(2, 2))
plot(model)
par(mfrow = c(1, 1))
influence.measures(model)
\end{lstlisting}

\begin{outputguidebox}
\textbf{Çözüm ve kontrol değerleri.} Model
$\widehat{\text{Kaygı}}=89{,}5403-4{,}9836\,\text{Uyku}$ biçimindedir.
Eğim standart hatası 1,6290, $t(18)=-3{,}061$ ve $p=0{,}00675$'tir.
$R^2=0{,}3421$ olur. Yedi saat için nokta tahmini yaklaşık 54,66'dır;
ortalama ve bireysel aralıklar ayrı \texttt{predict} çağrılarından okunur.
\end{outputguidebox}

\textbf{Yorum.} 
\texttt{interval="confidence"} aynı uyku süresindeki ortalama kaygı için,
\texttt{interval="prediction"} yeni bir öğrencinin puanı için aralık verir.
İkinci aralık bireysel değişkenliği içerdiğinden daha geniştir. Model tanıları
eğim anlamlı bulundu diye atlanmamalıdır.

\textbf{Pekiştirme ve yanıt.} Veri 4,5--8,5 saat aralığındadır. On iki saat
için öngörü hesaplanabilse bile güçlü bir dışa taşırmadır ve bilimsel olarak
ayrıca gerekçelendirilmelidir.
\section{Bölüm sonu çözümlü problemler}

\begin{enumerate}
  \item \textbf{Denklem ve tahmin.} $\widehat Y=70-2{,}5X$ modelinde $X=8$ için tahmini bulun ve eğimi yorumlayın.

  \textbf{Çözüm:} $\widehat Y=70-2{,}5(8)=50$'dir. X'teki bir birimlik artış, beklenen Y ortalamasında 2,5 birimlik azalmayla ilişkilidir. Gözlemsel veride bu ifade müdahale etkisi değildir.

  \item \textbf{Artık ve uyum.} Bir öğrencinin gözlenen puanı 58, model tahmini 52 ise artığı bulun. $SST=900$ ve $SSE=540$ için $R^2$ değerini hesaplayın.

  \textbf{Çözüm:} Artık $e=58-52=6$'dır; model öğrencinin puanını 6 birim düşük tahmin etmiştir. $R^2=1-SSE/SST=1-540/900=0{,}40$ olur.

  \item \textbf{İki aralığı ayırma.} $X=6$ için ortalama yanıt güven aralığı $[45;51]$, bireysel tahmin aralığı $[32;64]$ bulunmuştur. Farkı açıklayın.

  \textbf{Çözüm:} İlk aralık, $X=6$ olan birimlerin evren ortalamasını hedefler. İkinci aralık yeni tek bir bireyin doğal artık değişkenliğini de içerir; bu nedenle daha geniştir. Tek öğrencinin puanı ortalama güven aralığıyla değerlendirilmemelidir.
\end{enumerate}

\bolumkaynaklari{\cite{kutner2005,james2021,hastie2009,anscombe1973}}

\section*{Hatırla--Uygula--Yansıt}

\begin{enumerate}
\renewcommand{\labelenumi}{\textbf{\Alph{enumi}.}}
  \item \textbf{Hatırla:} $Y_i$, $X_i$, $\beta_0$, $\beta_1$ ve $\varepsilon_i$ terimlerini bir PDR örneği üzerinde açıklayın.
  \item \textbf{Yorumla:} $\widehat Y=72-3{,}5X$ denkleminde eğimi X ve Y'nin birimleriyle yorumlayın.
  \item \textbf{Kesme:} Aynı denklemde $X=0$ veri aralığının dışındaysa kesme yorumunun neden sınırlı olduğunu açıklayın.
  \item \textbf{Hesapla:} $b_1=-2{,}4$ ve $SE(b_1)=0{,}8$ için $t$ istatistiğini bulun.
  \item \textbf{Aralık:} Eğim için yüzde 95 güven aralığı $[-4{,}1;-0{,}7]$ olduğunda yönü ve makul büyüklükleri açıklayın.
  \item \textbf{Artık:} Gözlenen Y 60 ve uydurulan Y 54 ise artığı hesaplayıp yorumlayın.
  \item \textbf{Uyum:} $SST=1000$ ve $SSE=650$ için $R^2$ değerini bulun; ne söylediğini ve söylemediğini yazın.
  \item \textbf{Bağlantı:} Basit regresyonda $r=-0{,}50$ ise $R^2$ ve eğimin işareti nedir?
  \item \textbf{Aralıkları ayır:} Aynı X değeri için bireysel tahmin aralığının neden ortalama güven aralığından geniş olduğunu açıklayın.
  \item \textbf{Tanıla:} Artık--uydurulan değer grafiğinde eğri ve huni desenlerinin gösterebileceği sorunları ayırın.
  \item \textbf{Etki:} Büyük artık, yüksek kaldıraç ve yüksek etki kavramlarına ayrı örnekler verin.
  \item \textbf{Ekstrapolasyon:} 10--18 yaş verisinden 35 yaş için tahmin yapmanın riskini açıklayın.
  \item \textbf{Nedensellik:} Uyku--kaygı eğimi için ters nedensellik ve iki ortak neden önerin.
  \item \textbf{Uygula:} Python verisine $(9,90)$ gözlemini ekleyerek eğim, $R^2$, Cook uzaklığı ve grafikleri karşılaştırın.
  \item \textbf{Raporla:} Denklem, eğim, standart hata, test, güven aralığı, $R^2$, artık hatası ve nedensellik sınırı içeren kısa bir bulgu paragrafı yazın.
\end{enumerate}

\bolumbaglantisi{chapters/pdr/11-basit-regresyon.tex}